\documentclass[12pt]{article}
\usepackage[letterpaper,margin=0.75in]{geometry}
\usepackage[T1]{fontenc}
\usepackage{lmodern}
\usepackage{microtype}
\usepackage{booktabs}
\usepackage{enumitem}
\usepackage{needspace}
\usepackage{hyperref}
\usepackage{setspace}
\setstretch{1.0}
\setlength{\parindent}{1em}
\setlength{\parskip}{0pt}
\setlength{\tabcolsep}{5pt}
\setlist{nosep,leftmargin=1.5em}
\hypersetup{
  hidelinks,
  pdftitle={Cross-fitted causal erasure separates subject and attractor number in DistilGPT-2},
  pdfauthor={},
  pdfsubject={Anonymous PLC 51 abstract},
  pdfkeywords={}
}
\pagestyle{plain}
\newcommand{\ex}[1]{\textit{#1}}

\begin{document}

\begin{center}
{\bfseries Cross-fitted causal erasure separates subject and attractor\\
number in DistilGPT-2}
\end{center}
\vspace{-0.65em}

Agreement attraction places two number values in competition. In (1), the
hierarchical subject is singular and the closest noun is plural. This contrast
also exposes a familiar problem in mechanistic interpretation: a probe can
recover a feature that downstream layers ignore (Elazar et al. 2021), and a
specialized attention head may have little control over the output (Timkey \&
Linzen 2026). We test which number representation contributes to the next-word
preference of DistilGPT-2. Nearest-noun number reaches perfect probe accuracy,
while cross-fitted erasure of a subject-correlated direction produces a
1.578-logit loss relative to attractor-direction erasure. The effect appears
late and survives a lexical-cluster bootstrap.

\begin{enumerate}[label=(\arabic*),leftmargin=2.2em,topsep=2pt]
\item \ex{The key near the cabinets \{is/are\}\ldots}
\end{enumerate}
\vspace{-0.1em}

\textbf{Question.}
Causal work has localized agreement information to model components and
residual-stream directions (Finlayson et al. 2021; Ferrando \& Costa-juss\`a
2024). GPT-2 also shows attraction effects compatible with similarity-based
interference (Ryu \& Lewis 2021). These findings leave open whether the
accessible representation of the nearby noun contributes to inflection.
A nearest-noun account predicts a substantial effect when its fitted direction
is removed, with the largest cost on conflicting-number items. A hierarchical
account predicts a larger cost for the fitted subject direction.

\textbf{Materials and behavioral result.}
We generated 80 prefixes from ten lexical sets. The design crossed subject and
attractor number within two structures. In the PP
condition, the attractor was inside a subject-modifying PP as in (1); in the
relative-clause condition it was the embedded subject, as in \ex{The key that
the cabinets can describe}. All target forms, \ex{is} and \ex{are}, are single
tokens. We measured their next-token logit contrast in DistilGPT-2
(81.9M parameters). The model selected the grammatical form on 92.5\% of items
(74/80). Accuracy was 97.5\% for number-congruent items but 87.5\% for
incongruent items; it was 100\% in PPs and 85\% in relative clauses. The six
errors were concentrated in relative clauses, especially those with plural
subjects (75\% accuracy). We therefore treat the analysis as a controlled pilot.
Its scope does not extend to agreement behavior across the model family.

\textbf{Diagnostic analysis.}
At the final prefix position, we extracted the embedding state and the residual
stream after each of six Transformer blocks. Separate logistic probes predicted
subject and attractor number. Each of five folds held out two complete lexical
sets, including both structures and every number condition. Subject-number
accuracy is at chance in the embedding state and reaches 93.75\% after block 5.
Attractor number is easier to decode and reaches 100\% after blocks 3 and 4.

\begin{table}[h]
\centering
\caption{Group-cross-validated probe accuracy (\%) and cross-fitted loss in
the correct-inflection logit margin. State 0 is the embedding output.}
\vspace{-0.35em}
\begin{tabular}{@{}rrrrrr@{}}
\toprule
State & Subj. probe & Attr. probe & Subj. erase & Attr. erase & Random\\
\midrule
0 & 50.00 & 75.00 & -- & -- & --\\
1 & 71.25 & 86.25 & 0.003 & 0.001 & 0.000\\
2 & 58.75 & 97.50 & 0.008 & $-0.006$ & 0.002\\
3 & 73.75 & 100.00 & 0.019 & $-0.030$ & $-0.001$\\
4 & 91.25 & 100.00 & 0.332 & $-0.005$ & 0.000\\
5 & 93.75 & 97.50 & 0.712 & $-0.003$ & $-0.000$\\
6 & 92.50 & 88.75 & 1.578 & 0.000 & 0.003\\
\bottomrule
\end{tabular}
\end{table}
\vspace{-0.55em}

\textbf{Causal intervention.}
Each test item received a direction and center fit on the other eight lexical
sets. We separately projected its final-position state away from the fitted
subject and attractor directions, then ran the remaining blocks. Ten random
directions per fold supplied a dimensionality-matched control. The
subject-direction loss is small through block 3, then grows to 0.332, 0.712,
and 1.578 logits after blocks 4--6. Final accuracy becomes 66.25\%, against a
92.5\% baseline; incongruent accuracy becomes 55\%, against 87.5\%. The fitted attractor
direction changes the margin by less than 0.001 logits and leaves both
accuracies unchanged. Random-direction erasure changes the final margin by
0.003 logits. The final subject-minus-attractor contrast is 1.578 logits; a
10,000-sample bootstrap over lexical sets gives a 95\% CI of [1.460, 1.686].

\Needspace{4\baselineskip}
\textbf{Analysis.}
The residual stream preserves linearly recoverable information about the
controller and the nearby distractor. Their fitted directions have different
causal effects. The subject-correlated direction contributes to the
model's late logit contrast even at states where attractor number has higher
probe accuracy. Probe accessibility therefore overstates the causal role of
the fitted attractor direction.

The late onset agrees with Ferrando and Costa-juss\`a's (2024) subject-number
direction in the final residual stream and with Finlayson et al.'s (2021)
structure-dependent agreement mechanisms. The evidence supports a narrow
claim: in this pilot, a cross-fitted subject-correlated direction contributes
to the copular logit contrast. Under the same intervention, the most accessible
linear attractor direction has no measurable effect. The erasure effect is
similar across PPs and relative clauses (1.616 versus 1.540 logits), though it
is asymmetric by subject number. Final erasure lowers singular accuracy to
42.5\%; plural accuracy remains at 90\%. The fitted axis may combine number with
a singular-specific output bias. A one-dimensional null effect also leaves
other attractor representations untested. Replication should use larger
stimulus sets and model families, then compare projection methods such as LEACE
or mean projection (Dobrzeniecka et al. 2025).

\vspace{-0.35em}
\begin{thebibliography}{9}
\setlength{\itemsep}{0pt}
\setlength{\parskip}{0pt}

\bibitem[Dobrzeniecka et al. 2025]{dobrzeniecka}
Dobrzeniecka, Alicja, Antske Fokkens \& Pia Sommerauer. 2025.
Improving causal interventions in amnesic probing with mean projection or LEACE.
\textit{Findings of ACL}, 12981--12993.

\bibitem[Elazar et al. 2021]{elazar}
Elazar, Yanai, Shauli Ravfogel, Alon Jacovi \& Yoav Goldberg. 2021.
Amnesic probing: Behavioral explanation with amnesic counterfactuals.
\textit{TACL} 9. 160--175.

\bibitem[Ferrando \& Costa-juss\`a 2024]{ferrando}
Ferrando, Javier \& Marta R. Costa-juss\`a. 2024.
On the similarity of circuits across languages: A case study on subject-verb agreement.
\textit{Findings of EMNLP}, 10115--10125.

\bibitem[Finlayson et al. 2021]{finlayson}
Finlayson, Matthew, Aaron Mueller, Sebastian Gehrmann, Stuart Shieber,
Tal Linzen \& Yonatan Belinkov. 2021.
Causal analysis of syntactic agreement mechanisms in neural language models.
\textit{ACL-IJCNLP}, 1828--1843.

\bibitem[Ryu \& Lewis 2021]{ryu}
Ryu, Soo Hyun \& Richard Lewis. 2021.
Accounting for agreement phenomena in sentence comprehension with Transformer
language models. \textit{CMCL}, 61--71.

\bibitem[Timkey \& Linzen 2026]{timkey}
Timkey, William \& Tal Linzen. 2026.
Do syntax heads influence model behavior?
\textit{arXiv:2610.19231}.

\end{thebibliography}

\end{document}

